% Minimal retained preamble for lane-p2 cards taken from
% sources/logic-fl05-ramsey-like-theorems-for-the-schre/pinned-original-source.tex
% (arXiv:2412.11598v3, CC BY 4.0).  This source ships no separate preamble file:
% the pinned tex is a concatenation of the authors' main.tex and macros.tex and
% carries its preamble inline, so the blocks below are sliced mechanically from
% that file and nothing has been renamed or redefined.
%
% Copied verbatim from the source:
%   line 2        : the authors' class line
%   lines 3-5,7,9 : graphicx, amsthm, the ams bundle, cleveref, url
%   line 10       : tikz-cd, used by the authors' commutative diagrams
%   lines 20-108  : the authors' macro definitions, taken in full so that every
%                   macro of the retained spans is available
%   lines 113-124 : the authors' theorem declarations
%   lines 126-128 : the authors' macros for the three calligraphic letters, which
%                   sit just after the theorem declarations in the source
%   lines 131-141 : the authors' machinery for restating a theorem
%   lines 143-148 : the authors' coloured comment macros
%
% Omitted source preamble lines and why:
%   line 1, 110-111 and 150 (concatenation comment, title, date, author): the
%       assembler supplies the title block, author list and date;
%   line 6 (hyperref) and line 8 (inputenc with the utf8x option): the forced
%       xelatex engine supplies the hyperref wiring itself and utf8x is
%       incompatible with it;
%   line 115 (numberwithin theorem by section): the assembler numbers the
%       retained material itself;
%   lines 120 and 121-123 are declarations that the authors left commented out or
%       that duplicate earlier ones and are not reproduced.
\documentclass{article}
\usepackage{graphicx}
\usepackage{amsthm}
\usepackage{amsfonts, amsmath, amssymb, amsthm}
\usepackage{cleveref}
\usepackage{url}
\usepackage{tikz-cd}
\newcommand{\Psf}{\mathsf{P}}
\newcommand{\Qsf}{\mathsf{Q}}
\newcommand{\Rsf}{\mathsf{R}}
\def\F{\mathcal{F}}
\def\U{\mathcal{U}}
\def\C{\mathcal{C}}
\def\I{\mathcal{I}}
\def\D{\mathcal{D}}
\def\L{\mathcal{L}}
\def\V{\mathcal{V}}
\def\A{\mathcal{A}}
\def\B{\mathcal{B}}
\def\P{\mathcal{P}}
\def\M{\mathcal{M}}
\def\Nc{\mathcal{N}}
\def\Q{\mathcal{Q}}
\newcommand{\cs}{2^{\NN}}
\newcommand{\bs}{\NN^{\NN}}
\newcommand{\bstr}{2^{<\NN}}
\newcommand{\uh}{\upharpoonright}
\newcommand{\converge}{\!\!\downarrow}
\newcommand{\diverge}{\!\!\uparrow}
\newcommand{\qvdash}{\operatorname{{?}{\vdash}}}
\newcommand{\nqvdash}{\operatorname{{?}{\nvdash}}}
\newcommand{\ISig}{\mathsf{I}\Sigma^0}
\newcommand{\Robinson}{\mathsf{Q}}
\newcommand{\IPi}{\mathsf{I}\Pi^0}
\newcommand{\IDelta}{\mathsf{I}\Delta^0}
\newcommand{\Iopen}{\mathsf{I_{open}}}
\newcommand{\BSig}{\mathsf{B}\Sigma^0}
\newcommand{\BPi}{\mathsf{B}\Pi^0}
\newcommand{\BDelta}{\mathsf{B}\Delta^0}
\newcommand{\BCSig}{\mathsf{BC}\Sigma^0}
\newcommand{\BCPi}{\mathsf{BC}\Pi^0}
\newcommand{\BCDelta}{\mathsf{BC}\Delta^0}
\newcommand{\LSig}{\mathsf{L}\Sigma^0}
\newcommand{\LPi}{\mathsf{L}\Pi^0}
\newcommand{\CSig}{\mathsf{C}\Sigma}
\newcommand{\RCA}[0]{\mathsf{RCA}}
\newcommand{\WKL}[0]{\mathsf{WKL}}
\newcommand{\WWKL}[0]{\mathsf{WWKL}}
\newcommand{\ACA}[0]{\mathsf{ACA}}
\newcommand{\ATR}[0]{\mathsf{ATR}}
\newcommand{\PIOOCA}[0]{\Pi^1_1\mathsf{-CA}}
\newcommand{\CAC}[0]{\mathsf{CAC}}
\newcommand{\ADS}[0]{\mathsf{ADS}}
\newcommand{\RRT}[0]{\mathsf{RRT}}
\newcommand{\FS}[0]{\mathsf{FS}}
\newcommand{\PFS}[0]{\mathsf{PFS}}
\newcommand{\TS}[0]{\mathsf{TS}}
\newcommand{\GRRT}[0]{\mathsf{GRRT}}
\newcommand{\LRRT}[0]{\mathsf{LRRT}}
\newcommand{\EM}[0]{\mathsf{EM}}
\newcommand{\SEM}[0]{\mathsf{SEM}}
\newcommand{\HalfEM}[0]{\mathsf{Half}\mbox{-}\mathsf{EM}}
\newcommand{\RT}[0]{\mathsf{RT}}
\newcommand{\SRT}[0]{\mathsf{SRT}}
\newcommand{\NN}[0]{\mathbb{N}}
\newcommand{\nat}[0]{\mathbb{N}}
\newcommand{\card}{\operatorname{card}}
\newcommand{\finsub}{\subseteq_{\mathtt{fin}}}
\newcommand{\W}{\mathrm{W}}
\newcommand{\sW}{\mathrm{sW}}
\newcommand{\lex}{\mathrm{lex}}
\newcommand{\base}{\operatorname{base}}
\newcommand{\range}{\operatorname{ran}}
\def\qt#1{``#1''}%

\newtheorem*{statement}{Statement}
\newtheorem{theorem}{Theorem}
\newtheorem{maintheorem}[theorem]{Main Theorem}
\newtheorem{lemma}[theorem]{Lemma}
\newtheorem{question}[theorem]{Question}
\newtheorem{proposition}[theorem]{Proposition}
%\newtheorem{statement}[theorem]{Statement}
\newtheorem{remark}[theorem]{Remark}
\newtheorem{definition}[theorem]{Definition}
\newtheorem{corollary}[theorem]{Corollary}
\newtheorem{example}[theorem]{Example}
\newcommand{\R}{\mathcal{R}}
\renewcommand{\S}{\mathcal{S}}
\newcommand{\T}{\mathcal{T}}

% Allow to repeat theorems
\makeatletter
\newtheorem*{rep@theorem}{\rep@title}
\newcommand{\newreptheorem}[2]{%
\newenvironment{rep#1}[1]{%
 \def\rep@title{#2 \ref{##1}}%
 \begin{rep@theorem}}%
 {\end{rep@theorem}}}
\makeatother

\newreptheorem{theorem}{Theorem}

% Comments
\usepackage{xcolor}
\newcommand{\ludovic}[1]{\textcolor{blue}{#1}}
\newcommand{\quentin}[1]{\textcolor{green}{#1}}
\newcommand{\lorenzo}[1]{\textcolor{orange}{#1}}
\newcommand{\oriola}[1]{\textcolor{red}{#1}}
\title{The forcing question for a sigma-0-1 formula is equivalent to a first-order statement about one level of bounded progressive functions}
\author{Lorenzo Carlucci, Oriola Gjetaj, Quentin Le Houérou, Ludovic Levy Patey}
\date{arxiv:2412.11598v3; CC BY 4.0}
\begin{document}
\maketitle
\noindent\textit{Editorial note.} The counted unit is the article's own Lemma lem:fs-forcing-question-sigma01 at source lines 1201-1204 together with its proof at source lines 1205-1210, a straight adjacent pairing in the pinned file. The pinned file is a concatenation of the authors' main and macro files; it is the article's fixed version arXiv:2412.11598v3 under CC BY 4.0.\par\medskip
\noindent\textit{Editorial note.} This source ships no separate preamble file, so the retained preamble is sliced mechanically from the pinned file by the lane and is shared, byte for byte, with the other cards of this source; it carries the authors' full macro block, including the symbols for the forcing relation, for the tree and for the space of functions that the retained statement and proof use.\par\medskip
\noindent\textit{Editorial note.}  The retained preamble is the lane's shared slice for this source; it was extended during this card's first revision to carry the three macros for calligraphic R, S and T, which in the source sit just after the theorem declarations and which the statement of this lemma needs. The first revision failed on exactly that undefined macro, and this revision is the corrected one.\par\medskip
\noindent\textit{Editorial note.} One context span carries the definition of the forcing question at source lines 1194-1197, which the retained statement names; it is context, not a counted unit, and it does not overlap the counted interval.\par\medskip
\noindent\textit{Editorial note.} The retained statement and proof contain no reference to another labelled result of the article and no citation: the only external ingredient is the standard form of Konig's lemma applied to the finitely branching tree of counterexamples, which the author names in passing and which is recorded in the standard prerequisites. Accordingly this card carries no navigation pointer and no bibliography entry.\par\medskip
\noindent\textit{Editorial note.} The article states immediately above the counted lemma that the forcing question as defined is of second-order complexity, and the counted lemma is its first-order reformulation; that surrounding prose is not part of the counted unit and is not retained.\par\medskip
\noindent\textit{Allowed standard prerequisites (editorial).}\begin{itemize}
\item The forcing question, the space of relations and its correspondence with bounded functions. The section fixes a condition, a formula, and a relation space whose members are relations between finite sets of natural numbers and natural numbers, bounded by a computable function, and in one-to-one correspondence with the constrained functions: to a function one attaches the relation of the pairs lying in its values, and conversely. The forcing question holds for a condition and a formula when for every relation of that space there is a finite block inside the reservoir that is free for the corresponding function and makes the formula hold on the enlarged stem; the formula considered here is of the lowest syntactic complexity, so its truth is witnessed by a finite amount of information. The author's own definition and the two lemmas of the section about the forcing question and its specification are the allowed inputs.. The counted lemma is the step that turns the second-order formulation of the forcing question into a first-order statement about one level of bounded progressive functions, which is what makes the generic filter construction effective.
\end{itemize}\medskip
\begin{definition}
Let $(f, \sigma, X)$ be a condition and $\varphi(G)$ be a $\Sigma^0_1$-formula.
Let $(f, \sigma, X) \qvdash \varphi(G)$ hold if and only if for every relation $R \in \R_{b_f^{+}}$, there is a finite $g_R$-free set $\rho \subseteq X$ such that $\varphi(\sigma \cup \rho)$ holds.
\end{definition}

\begin{lemma}\label[lemma]{lem:fs-forcing-question-sigma01}
Let $(f, \sigma, X)$ be a condition and $\varphi(G)$ be a $\Sigma^0_1$-formula.
Then $(f, \sigma, X) \qvdash \varphi(G)$ if and only if there is some~$k \in \omega$ such that for every $b_f^{+}$-constrained progressive function $g : [0, k]^{\leq h(\cdot)} \to [0, k]^{<\omega}$, there is a finite $g$-free set $\rho \subseteq X \uh_k$ such that $\varphi(\sigma \cup \rho)$ holds.
\end{lemma}

\begin{proof}
Suppose first that there is some~$k \in \NN$ such that for every $b_f^{+}$-constrained progressive function $g : [0, k]^{\leq h(\cdot)} \to [0, k]^{<\omega}$, there is a finite $g$-free set $\rho \subseteq X \uh_k$ such that $\varphi(\sigma \cup \rho)$ holds. Let $R \in \R_{b_f^+}$ be a relation,
and let $g_R \in \P_{b_f^+}$ be the corresponding function. Then, letting $g : [0, k]^{\leq h(\cdot)} \to [0, k]^{<\omega}$ be defined by $g(s) = g_R(s) \cap [0, k]$, the function $g$ is $b_f^{+}$-constrained and progressive, so there is some finite $g$-free set $\rho \subseteq X \uh_k$ such that $\varphi(\sigma \cup \rho)$ holds. In particular, $\rho$ is $g_R$-free. Since there is such a $\rho$ for every $R \in \R_{b_f^+}$, then $(f, \sigma, X) \qvdash \varphi(G)$ holds.

Suppose now that for every~$k \in \NN$, there is a $b_f^{+}$-constrained progressive function $g_k : [0, k]^{\leq h(\cdot)} \to [0, k]^{<\omega}$ such that for every $g_k$-free set $\rho \subseteq X \uh_k$, $\varphi(\sigma \cup \rho)$ does not hold. Let $\T$ be the tree which, at level~$k$, contains all such functions $g_k$, and which is ordered by the function extension relation. The tree $\T$ is finitely branching, so by K\"onig's lemma, there is an infinite path $g \in \T$. This path is a function $g \in \P_{b_f^+}$ such that for every finite $g$-free set $\rho \subseteq X$, $\varphi(\sigma \cup \rho)$ does not hold. Then the relation $R_g$ witnesses that $(f, \sigma, X) \qvdash \varphi(G)$ does not hold.
\end{proof}

\end{document}
